% Source: https://artofproblemsolving.com/community/c3h2037908_divisible_with_7
% Stable identity: aops-topic-2037908-q01
Sort the \(37\) numbers into the \(7\) classes according to their remainders \(0,1,\dots,6\) modulo \(7\).

\emph{Case 1: some class contains at least \(7\) of the numbers.} Choose \(7\) numbers from it; all have the same remainder \(r\), so their sum is \(\equiv7r\equiv0\pmod 7\).

\emph{Case 2: every class contains at most \(6\) numbers.} If some class were empty, the \(37\) numbers would lie in at most \(6\) classes with at most \(6\) each, i.e. at most \(36\) numbers, a contradiction. So every class is nonempty. Choose one number from each class; their sum is
\[
\equiv0+1+2+3+4+5+6=21\equiv0\pmod 7 .
\]
In both cases we find \(7\) numbers whose sum is divisible by \(7\). \(\blacksquare\)
